\documentclass[../main.tex]{subfiles}
\begin{document}

% ============================================================
%  §4.1 Μοντέλο Οχήματος -- δίτροχο μοντέλο, συντεταγμένες, weight transfer
%  Ενσωματωμένα TikZ σχήματα inline.
% ============================================================

\label{sec:vehiclemodel}

Η ενσωμάτωση του μοντέλου ελαστικών σε έναν προσομοιωτή γύρου απαιτεί τον ορισμό ενός μοντέλου οχήματος που να συνδέει τα κινηματικά μεγέθη (ταχύτητα, επιτάχυνση) με τις δυνάμεις που αναπτύσσονται σε κάθε τροχό. Στην παρούσα εργασία χρησιμοποιείται ένα \emph{δίτροχο μοντέλο} (two-track model) τριών βαθμών ελευθερίας: διαμήκης επιτάχυνση, πλευρική επιτάχυνση και στρέψη (yaw). Το μοντέλο αντιμετωπίζει το σώμα ως άκαμπτο σώμα (rigid body) χωρίς δυναμική ανάρτησης -- οι κατακόρυφες μετατοπίσεις μοντελοποιούνται μόνο μέσω της στατικής κατανομής φορτίου.

Η επιλογή δίτροχου έναντι μονότροχου (bicycle model) μοντέλου είναι εδώ αναγκαία: στο δίτροχο διακρίνονται οι αριστεροί από τους δεξιούς τροχούς, γεγονός κρίσιμο για τον υπολογισμό της πλευρικής μεταφοράς φορτίου και για την ανεξάρτητη παρακολούθηση της θερμοκρασίας κάθε τροχού.

% ---
\subsection{Συστήματα συντεταγμένων}
\label{subsec:frames}
% ---

Για τη μαθηματική περιγραφή της κίνησης χρησιμοποιούνται τρία αλληλοεξαρτώμενα συστήματα συντεταγμένων, τα οποία ορίζονται σύμφωνα με τις συμβάσεις SAE \cite{Milliken1995}. Στο \cref{fig:twotrack_topview} φαίνεται η γεωμετρική σχέση τους.

% ----- TikZ: κάτοψη δίτροχου μοντέλου -----
\begin{figure}[htbp]
\centering
\begin{tikzpicture}[scale=1.05,
    wheel/.style={rectangle, draw=black, fill=gray!30,
                  minimum width=0.28cm, minimum height=0.72cm, rounded corners=1pt},
    frameaxis/.style={->, thick},
    carframe/.style={->, thick, red!70!black},
    tyreframe/.style={->, thick, green!50!black},
    dim/.style={<->, thin, gray}
]
  % -- car body --
  \draw[fill=gray!15, draw=black, rounded corners=3pt] (-1.0,-0.55) rectangle (1.0,0.55);
  \node[font=\scriptsize] at (0,0) {Σώμα / CG};

  % -- wheels --
  \node[wheel] (FL) at (-0.65, 0.85) {};  \node[font=\scriptsize,above=1pt] at (FL.north) {FL};
  \node[wheel] (FR) at (-0.65,-0.85) {};  \node[font=\scriptsize,below=1pt] at (FR.south) {FR};
  \node[wheel] (RL) at ( 0.65, 0.85) {};  \node[font=\scriptsize,above=1pt] at (RL.north) {RL};
  \node[wheel] (RR) at ( 0.65,-0.85) {};  \node[font=\scriptsize,below=1pt] at (RR.south) {RR};

  % -- global frame (bottom-left) --
  \begin{scope}[shift={(-2.8,-1.7)}]
    \draw[frameaxis,black] (0,0)--(0.85,0) node[right,font=\small]{$X$};
    \draw[frameaxis,black] (0,0)--(0,0.85) node[above,font=\small]{$Y$};
    \node[font=\scriptsize,below left] at (0,0) {Global};
  \end{scope}

  % -- car frame (rotated β=10°) --
  \begin{scope}[rotate=10]
    \draw[carframe] (0,0)--(0.72,0) node[right,font=\scriptsize,red!70!black]{$x$};
    \draw[carframe] (0,0)--(0,0.72) node[above,font=\scriptsize,red!70!black]{$y$};
  \end{scope}
  \node[font=\scriptsize,red!70!black] at (0.55,0.55) {Car frame};
  % beta arc
  \draw[->,thin,purple] (0.42,0) arc[start angle=0,end angle=10,radius=0.42];
  \node[font=\scriptsize,purple] at (0.68,0.12) {$\beta$};

  % -- tyre frame front (δ=20°) --
  \begin{scope}[shift={(-0.65,0.85)},rotate=20]
    \draw[tyreframe] (0,0)--(0.35,0);
  \end{scope}
  \draw[->,thin,green!50!black] (-0.65+0.22,0.85) arc[start angle=0,end angle=20,radius=0.22];
  \node[font=\scriptsize,green!50!black] at (-0.15,1.12) {$\delta$};

  % -- velocity vector --
  \draw[->,ultra thick,orange!80!black] (0,0)--(-0.55,0.22)
    node[left,font=\small,orange!80!black]{$\mathbf{v}$};

  % -- dimensions --
  \draw[dim] (-0.65,1.2)--(0.65,1.2);
  \node[gray,font=\scriptsize,above] at (0,1.2) {$L$};
  \draw[dim] (-1.1,0.85)--(-1.1,-0.85);
  \node[gray,font=\scriptsize,left] at (-1.1,0) {$t$};

  % -- CG dot --
  \filldraw[black] (0,0) circle (2pt);
\end{tikzpicture}
\caption{Κάτοψη δίτροχου μοντέλου. Αδρανειακό σύστημα (μαύρο), σύστημα σώματος (κόκκινο) υπό γωνία πλαϊνής ολίσθησης $\beta$, γωνία διεύθυνσης εμπρός τροχού $\delta$ (πράσινο) και διάνυσμα ταχύτητας $\mathbf{v}$ (πορτοκαλί).}
\label{fig:twotrack_topview}
\end{figure}
% ----- τέλος TikZ -----

Το \emph{αδρανειακό (global) σύστημα} έχει άξονα $X$ προς τα εμπρός και $Y$ προς τα αριστερά. Στο πλαίσιο αυτό ορίζονται οι μετρούμενες επιταχύνσεις $a_x$ (διαμήκης) και $a_y$ (πλευρική) και ισχύει η εξίσωση ισορροπίας $m\mathbf{a} = \sum_i \mathbf{F}_i$.

Το \emph{σύστημα σώματος (car frame)} είναι σταθερό ως προς το αμάξωμα και περιστρέφεται από το global κατά τη γωνία πλαϊνής ολίσθησης $\beta$ (chassis side-slip angle). Η γωνία $\beta$ εκφράζει την απόκλιση της κατεύθυνσης κίνησης του CG από τον άξονα του οχήματος -- στη συνήθη οδήγηση είναι μικρή ($<5°$), αλλά σε έντονες μανούβρες μπορεί να φτάσει έως $\pm 15°$.

Το \emph{σύστημα τύπου (tyre frame)} χρησιμοποιείται ξεχωριστά για κάθε τροχό. Για τους πίσω τροχούς ταυτίζεται με το car frame. Για τους εμπρός τροχούς περιστρέφεται επιπλέον κατά τη γωνία διεύθυνσης $\delta$ -- στο πλαίσιο αυτό ορίζονται οι δυνάμεις $F_x$ (διαμήκης) και $F_y$ (πλευρική) που αναπτύσσει κάθε τροχός.

% ---
\subsection{Κατανομή κατακόρυφων φορτίων -- Μεταφορά φορτίου}
\label{subsec:weight_transfer}
% ---

Ο κατακόρυφος φόρτος κάθε τροχού καθορίζει άμεσα το μέγιστο εφαπτομενικό που μπορεί να αναπτυχθεί (μέσω του $\mu \cdot F_z$) και την αναλογία θερμότητας που παράγεται ($Q_2 \propto F_z$). Η κατανομή των κατακόρυφων φορτίων αλλάζει δυναμικά με τις επιταχύνσεις.

% ----- TikZ: weight transfer -----
\begin{figure}[htbp]
\centering
\begin{tikzpicture}[scale=0.88,
    vec/.style={->,very thick},
    dim/.style={<->,thin,gray}
]
  % ground
  \draw[thick,gray] (-3.4,0)--(3.4,0);
  \fill[gray!20] (-3.4,-0.15) rectangle (3.4,0);

  % wheels (side view)
  \filldraw[gray!40,draw=black] (-2.1,0.28) circle (0.28);
  \filldraw[gray!40,draw=black] ( 2.1,0.28) circle (0.28);

  % body
  \draw[fill=gray!15,draw=black,rounded corners=3pt] (-2.3,0.56) rectangle (2.3,1.16);
  \node[font=\small] at (0,0.86) {Σώμα};

  % CG
  \filldraw[red] (-0.3,1.16) circle (3pt);
  \node[font=\small,above,red] at (-0.3,1.2) {CG};
  \draw[dashed,gray!60] (-0.3,0.56)--(-0.3,1.6);

  % hCG
  \draw[dim] (-3.0,0.56)--(-3.0,1.16);
  \node[gray,font=\small,left] at (-3.0,0.86) {$h_{\rm CG}$};

  % ax arrow
  \draw[vec,orange!80!black] (-0.3,1.7)--(1.3,1.7)
    node[right,font=\small]{$a_x>0$};

  % Fz arrows (static + delta)
  \draw[vec,blue!70] (-2.1,0.56)--(-2.1,-0.7) node[below,font=\small,blue!70]{$F_{z,f}^{\rm stat}$};
  \draw[vec,blue!70] ( 2.1,0.56)--( 2.1,-0.9) node[below,font=\small,blue!70]{$F_{z,r}^{\rm stat}$};

  % delta Fz
  \draw[vec,red!70,dashed] (-2.1,-0.45)--(-2.1,-0.1);
  \node[font=\scriptsize,red!70,left] at (-2.1,-0.28) {$-\Delta F_z$};
  \draw[vec,red!70,dashed] (2.1,-0.45)--(2.1,-0.9);
  \node[font=\scriptsize,red!70,right] at (2.1,-0.67) {$+\Delta F_z$};

  % wheelbase
  \draw[dim] (-2.1,-1.3)--(2.1,-1.3);
  \node[gray,font=\small,below] at (0,-1.3) {$L$};

  % equation
  \node[font=\small] at (0,-1.9)
    {$\Delta F_{z} = \dfrac{m\,a_x\,h_{\rm CG}}{L}$};
\end{tikzpicture}
\caption{Διαμήκης μεταφορά φορτίου κατά επιτάχυνση ($a_x>0$). Ο εμπρός άξονας αποφορτίζεται κατά $\Delta F_z$ και ο πίσω επιφορτίζεται αντίστοιχα.}
\label{fig:weight_transfer}
\end{figure}
% ----- τέλος TikZ -----

\paragraph{Διαμήκης μεταφορά:}
\begin{equation}
    \Delta F_{z,\text{long}} = \frac{m \cdot a_x \cdot h_{\text{CG}}}{L}
    \label{eq:Fz_long}
\end{equation}

Κατά επιτάχυνση ($a_x > 0$) το φορτίο μεταφέρεται από τον εμπρός στον πίσω άξονα. Κατά φρένα­ρισμα ($a_x < 0$) το αντίθετο -- ο εμπρός άξονας φορτίζεται περισσότερο, γι' αυτό και τα εμπρός ελαστικά υφίστανται εντονότερη φθορά φρεναρίσματος.

\paragraph{Πλευρική μεταφορά:}
\begin{equation}
    \Delta F_{z,\text{lat}} = \frac{m \cdot a_y \cdot h_{\text{CG}}}{t}
    \label{eq:Fz_lat}
\end{equation}

όπου $t$ ο τροχόδρομος (track width) του αντίστοιχου άξονα. Σε δεξιά στροφή ($a_y < 0$ στη σύμβαση $Y$-αριστερά) ο εξωτερικός (αριστερός) τροχός επιφορτίζεται. Αυτή η ασυμμετρία εξηγεί γιατί σε πίστες με κυριαρχία δεξιών στροφών (π.χ. Βαρκελώνη) τα αριστερά ελαστικά θερμαίνονται πιο έντονα.

% TODO: Αν χρησιμοποιείς roll stiffness distribution (anti-roll bars) για το
% πραγματικό split εμπρός/πίσω, αναφέρτο εδώ. Επίσης αξίζει σημείωση για
% aerodynamic downforce: αν το μοντέλο το συμπεριλαμβάνει ως συνάρτηση v².

% ---
\subsection{Επίλυση κατάστασης οχήματος}
\label{subsec:car_state_solver}
% ---

Σε κάθε σημείο $s_k$ της πίστας, γνωρίζουμε $v_k$, $a_{x,k}$ και $a_{y,k}$ (από το προφίλ ταχύτητας και ακτίνα καμπυλότητας, §\ref{subsec:slip_estimation}). Οι δύο άγνωστες ποσότητες που χαρακτηρίζουν την κατάσταση του οχήματος είναι η γωνία διεύθυνσης $\delta$ και η γωνία πλαϊνής ολίσθησης $\beta$.

Η ισορροπία δυνάμεων στο αδρανειακό πλαίσιο απαιτεί:
\begin{equation}
    \sum_{i \in \{FL,FR,RL,RR\}} \mathbf{F}_{i,\text{global}}(\delta,\beta) = m \cdot \mathbf{a}
    \label{eq:force_balance}
\end{equation}

Αυτό αποτελεί σύστημα 2 εξισώσεων (X-Y) με 2 αγνώστους, το οποίο επιλύεται αριθμητικά μέσω \texttt{lsqnonlin} (Newton-Raphson based) με φυσικά όρια $\delta \in [-45°, 45°]$, $\beta \in [-15°, 15°]$. Εσωτερικά, κάθε αξιολόγηση του residual περιλαμβάνει αλυσίδα υπολογισμών: μετασχηματισμοί συντεταγμένων $\to$ κατανομή $F_z$ $\to$ ταχύτητες τροχών $\to$ slip $\to$ δυνάμεις ελαστικού $\to$ global frame.

\end{document}
